\subsection{Hahn-Banach 定理及其应用}

定理 1（Hahn-Banach）. --- 设 V 是复向量空间 E 的一个向量子空间，f 是 V 上的一个（复）线性型，p 是 E 上的一个半范数，满足 \(|f(y)| \leqslant p(y)\)（对所有 \(y \in V\)）。则存在 E 上的一个线性型 \(f_{1}\)，它延拓 f 且满足 \(|f_{1}(x)| \leqslant p(x)\)（对所有 \(x \in E\)）。

因为 \(g = \mathcal{R}f\) 是定义在 \(V\) 中的实线性型，且满足 \(|g(y)| \leqslant p(y)\) 于 \(V\) 的每一点；因此存在一个实线性型 \(g_1\) 于 \(E\) 中，它延拓 \(g\) 且满足 \(|g_1(x)| \leqslant p(x)\)（对所有 \(x \in E\)）（II, p.~23, 推论 1）。设 \(f_1(x) = g_1(x) - ig_1(ix)\) 为 \(E\) 上以 \(g_1\) 为实部的那个复线性型（II, p.~61）。对所有实数 \(\vartheta\)

\[
\left| \mathcal {R} (e ^ {i \vartheta} f _ {1} (x)) \right| = \left| \mathcal {R} (f _ {1} (e ^ {i \vartheta} x)) \right| = \left| g _ {1} (e ^ {i \vartheta} x) \right| \leqslant p (e ^ {i \vartheta} x) = p (x)
\]

由于 \(p\) 是复空间 \(\mathbf{E}\) 上的一个半范数；这蕴含关系 \(|f_1(x)| \leqslant p(x)\)，定理得证。

推论 1. --- 设 \(x_0\) 是复拓扑向量空间 \(E\) 中的一点，\(p\) 是 \(E\) 中的一个连续半范数；则存在一个连续（复）线性型 \(f\)，它定义在 \(E\) 中，使得 \(f(x_0) = p(x_0)\)，且 \(|f(x)| \leqslant p(x)\)（对所有 \(x \in E\)）。

推论 2. --- 设 V 是复局部凸空间 E 的一个向量子空间，f 是 V 中定义且连续的一个（复）线性型；则存在 E 中定义且延拓 f 的一个连续线性型 \(f_{1}\)。若 E 是赋范的，则存在这样的线性型 \(f_{1}\)，它还满足 \(\|f_{1}\|=\|f\|\)。

推论 3. --- 设 M 是 Hausdorff 复局部凸空间 E 的一个有限维向量子空间。则存在 E 的一个闭向量子空间 N，它是 M 在 E 中的拓扑余子空间。

这些利用 p.~24 定理 1 的证明，与 II, p.~23 推论 2、p.~24 推论 3、p.~24 命题 2 以及 p.~25 推论 2 中的证明相同。

命题 1. --- 设 A 是复拓扑向量空间 E 中的一个非空开凸集，M 是一个不与 A 相交的非空（复）线性流形。则存在一个包含 M 且不与 A 相交的闭复超平面 H。

我们可以假设 \(0 \in M\)。于是存在一个包含 M 且不与 A 相交的闭实超平面 \(H_{0}\)（II, p.~36; 定理 1）。由于 M = iM，闭复超平面 \(\mathbf{H} = \mathbf{H}_{0} \cap (i\mathbf{H}_{0})\) 就有所要求的性质。

推论. --- 在复局部凸空间 E 中，每个闭复线性流形 M 都是包含它的那些闭复超平面的交。

事实上，对所有 \(x \notin M\)，存在 x 的一个不与 M 相交的开凸邻域 V，从而存在一个包含 M 且不与 V 相交的闭复超平面 H；更强的，H 不包含 x。

命题 2. --- 设 A 是复拓扑向量空间 E 的一个非空均衡开凸集，B 是一个不与 A 相交的非空凸集。则存在 E 上的一个连续复线性型 f 和一个数 \(\alpha > 0\)，使得 \(|f(x)| < \alpha\) 在 A 中成立，且 \(|f(y)| \geqslant \alpha\) 在 B 中成立。

因为，存在连续实线性型 \(g\) 于 \(\mathbf{E}\) 上，以及一个实数 \(\alpha\)，使得 \(g(x) < \alpha\) 在 \(\mathbf{A}\) 中成立，且 \(g(y) \geqslant \alpha\) 在 \(\mathbf{B}\) 中成立（II, p.~37, 命题 1）。由于 \(0 \in \mathbf{A}\)，我们有 \(\alpha > 0\)。我们证明，连续复线性型 \(f(x) = g(x) - ig(ix)\) 与数 \(\alpha\) 具有所要求的性质。因为，由于 \(\mathcal{R}f = g\)，我们有 \(|f(y)| \geqslant \alpha\) 在 \(\mathbf{B}\) 中成立。另一方面，对所有 \(x \in \mathbf{A}\) 与所有实数 \(\vartheta\)，点 \(e^{i\vartheta}x\) 属于 \(\mathbf{A}\)（因为 \(\mathbf{A}\) 是均衡的），且我们有 \(f(x) = e^{-i\vartheta}f(e^{i\vartheta}x)\)；于是存在一个数 \(\vartheta\)，使得 \(|f(x)| = \mathcal{R}(e^{i\vartheta}f(x)) = g(e^{i\vartheta}x) < \alpha\)，命题得证。

命题 3. --- 设 A 是复局部凸空间 E 中的一个均衡闭凸集，K 是 E 中一个不与 A 相交的非空紧凸集。则存在 E 上的一个连续复线性型 f 和一个数 \(\alpha > 0\)，使得 \(|f(x)| < \alpha\) 在 A 中成立，且 \(|f(y)| > \alpha\) 在 K 中成立。

本命题由 II, p.~38 命题 4 得出，正如命题 2 由 II, p.~37 命题 1 得出一样。

